\begin{abstract}
Assessing the fairness of a decision making system with respect to a protected
class, such as gender or race, is challenging when class membership labels
are unavailable.
%due to legal or operational reasons.
Probabilistic models for predicting the protected class based on observable proxies,
such as surname and geolocation for race, are sometimes used to
impute these missing labels for compliance assessments.
%purpose.
%Such methods are known to be used by regulators.
Empirically, these methods are observed to exaggerate disparities, but the
reason why is unknown.
In this paper, we decompose the biases in estimating outcome disparity via
%iexisting 
threshold-based imputation into multiple interpretable bias sources,
allowing us to explain when over- or underestimation occurs.
We also propose an alternative weighted estimator that uses soft classification,
%rules rather than hard imputation,
and show that its bias arises simply from the conditional covariance of
the outcome with the true class membership.
Finally, we illustrate our results with numerical simulations and a public dataset 
of mortgage applications, using geolocation as a proxy for race.
We confirm that the bias of threshold-based imputation is generally upward, but its
magnitude varies strongly with the threshold chosen.
%due to the complex interplay of multiple sources of bias.
Our new weighted estimator tends to have a negative bias that is much simpler
to analyze and reason about.
\end{abstract}
